% Plano da reta final da P1 em uma página (LuaLaTeX).
% Espelha prova/README.md §5 ("A reta final, hora a hora" e "Dúvidas para o professor"):
% mudou lá, muda aqui. Compilar: Rscript caderno/montar_caderno.R plano
\documentclass[10pt]{article}
\usepackage[a4paper,top=1.1cm,bottom=1.1cm,left=1.2cm,right=1.2cm]{geometry}
\usepackage{fontspec}
\usepackage{unicode-math}
\setmainfont{Libertinus Serif}
\setsansfont{Libertinus Sans}
\setmathfont{Libertinus Math}
\usepackage{polyglossia}
\setmainlanguage{brazil}
\usepackage{multicol}
\usepackage{tabularx}
\usepackage{booktabs}
\usepackage{enumitem}
\setlist{nosep,leftmargin=1.3em}
\setlength{\parindent}{0pt}
\setlength{\columnsep}{0.7cm}
\setlength{\columnseprule}{0.3pt}
\renewcommand{\arraystretch}{1.12}
\pagestyle{empty}

\newcommand{\dia}[1]{\par\medskip{\sffamily\bfseries\large #1}\par\smallskip}
\newcolumntype{H}{>{\raggedright\arraybackslash}p{2.15cm}}
\newcolumntype{Y}{>{\raggedright\arraybackslash}X}
\newenvironment{horas}{\tabularx{\linewidth}{@{}HY@{}}\toprule
  \textsf{\bfseries Hora} & \textsf{\bfseries O quê}\\\midrule}{\bottomrule\endtabularx}

\begin{document}
\fontsize{11pt}{13.4pt}\selectfont

{\sffamily\bfseries\LARGE Reta final da P1}\hfill{\sffamily seg 02/11, 20h30 \textrightarrow{} qui 05/11, prova às 10h}

\smallskip
\textbf{A lógica:} aprender pelo guia, testar no Simulado 02 no horário da prova, corrigir pelo guia, consolidar sem matéria nova.
\textbf{Em cada questão-tipo (QT):} (1) ler o enunciado e o “Antes de escrever”; (2) \textbf{cobrir o resto e tentar no papel}; (3) conferir linha a linha e, na primeira linha diferente, ler o sentido dela.

\begin{multicols}{2}\raggedright

\dia{Segunda 02/11 · 20h30–23h}
\begin{horas}
20h30–20h45 & básico, camadas 1 e 2, só leitura\\
20h45–21h30 & \textbf{QT4, ARMA(1,1)}: a que travou no S1\\
21h30–21h40 & pausa\\
21h40–22h10 & QT7 (AIC, BIC, MQO) e QT8 (previsão)\\
22h10–22h30 & QT9 (SARIMA)\\
22h30–22h50 & \textbf{QT4 de novo, folha branca, até 15 min}\\
22h50–23h & erros no log; dormir. \textbf{Não abrir o S2}\\
\end{horas}

\dia{Terça 03/11 · o dia todo}
\begin{horas}
8h30–8h50 & camadas 1 e 2 do básico \textbf{de memória}; conferir\\
8h50–9h50 & QT6, QT2 e QT3, cobrindo e tentando\\
9h50–10h & pausa\\
\textbf{10h–11h50} & \textbf{Simulado 02 em modo prova}: 3 min de leitura, 25 min de contas curtas, derivações, frase de fecho\\
12h–13h30 & almoço e descanso\\
13h30–14h30 & correção pelo gabarito, por passo; log; placar S2\\
14h30–16h & \textbf{cada item com perda}: a QT de mesmo número no guia, o sentido das linhas que faltaram, o item refeito em folha branca\\
16h–16h15 & pausa\\
16h15–17h15 & QT5 e QT1, cobrindo e tentando\\
17h15–18h & S3, \textbf{só as contas}: Q3c, Q4b–c, Q5a–b\\
18h–19h30 & jantar e descanso\\
19h30–20h30 & \textbf{QT4 de cabeça} (+1 dia) e S1 Q2 de cabeça\\
20h30–21h & as \textbf{9 frases de fecho} de memória; conferir\\
21h–21h30 & lista de dúvidas para o professor\\
21h30 & parar\\
\end{horas}

\columnbreak

\dia{Quarta 04/11 · noite (revista após o quadro)}
\begin{horas}
18h–19h & chegar, comer, descansar; nada de matéria\\
19h–19h20 & ler a folha \textbf{o que cai}\\
19h20–20h20 & \textbf{Simulado 05}: 4 × 2,5, em 1 hora\\
20h20–20h40 & correção pelo gabarito; log\\
20h40–21h20 & refazer cada item com perda, pelo guia e pela folha\\
21h20–21h50 & desenhar de memória FAC e FACP dos 6 padrões; frases de fecho\\
\textbf{até 22h30} & \textbf{parar e dormir}\\
\end{horas}

\dia{Quinta 05/11 · prova às 10h}
\textbf{7h} acordar. \textbf{8h}, 15 min: o básico, os quatro movimentos do ARMA e a ordem da prova.
\begin{enumerate}[label=\textbf{\arabic*.}]
\item 3 min de leitura: convenção do MA, $T$ e $s$;
\item cerca de 25 min \textbf{só de contas curtas};
\item derivações, da mais curta para a mais longa;
\item \textbf{frase de fecho} em cada item.
\end{enumerate}

\dia{O que o professor anunciou (04/11)}
\begin{itemize}
\item \textbf{4 questões de 2,5}: processos ($z_3=z_1+2z_2$, passeio); ARMA (estacionário, \textbf{montar a FAC}: conta e desenho); SARIMA (SMA$_4$ desenhado, equação de previsão); ciclo de modelagem e análise.
\item AIC, BIC, RMSE e MAPE: \textbf{interpretar}, não calcular.
\item Uma questão de \textbf{previsão}.
\end{itemize}

\dia{Fica de fora}
\textbf{Yule-Walker (riscado)}, Simulado 04, verossimilhança exata, testes de raiz unitária, Beveridge-Nelson, MASE, DM e CRPS, trabalho do INPC.

\end{multicols}

\vfill
\hrule\smallskip
{\footnotesize\raggedright\textbf{Onde está cada coisa}, em \texttt{caderno/pdf/}: o que cai (\texttt{caderno\_o\_que\_cai}) · Simulado 05 e gabarito (\texttt{caderno\_simulado\_05}) · guia das 9 QT (\texttt{caderno\_nao\_posso\_errar}) · básico (\texttt{caderno\_basico\_p1}) · tudo (\texttt{caderno\_p1\_completo}) · log em \texttt{prova/log\_erros.md}.\par}
\end{document}
